\documentclass[aspectratio=169]{beamer}
\usepackage{tikz}
\usepackage{graphicx}
\usepackage{listings}
\usepackage{xcolor}
\usepackage{hyperref}
\usepackage[spanish,es-noquoting]{babel}
\usepackage[round]{natbib}

\lstset{
  basicstyle=\ttfamily\small,
  keywordstyle=\color{blue},
  stringstyle=\color{red},
  commentstyle=\color{gray},
  showstringspaces=false
}
% Minimal look
\usetheme{Warsaw}
%\setbeamertemplate{headline}
\setbeamertemplate{footline}{
  \leavevmode%
  \hbox{%
    % Left: Section name
    \begin{beamercolorbox}[wd=.33\paperwidth,ht=2.25ex,dp=1ex,left]{author in head/foot}%
      \hspace*{2ex}\insertsection
    \end{beamercolorbox}%
    % Center: Presentation title
    \begin{beamercolorbox}[wd=.34\paperwidth,ht=2.25ex,dp=1ex,center]{title in head/foot}%
      \insertshorttitle
    \end{beamercolorbox}%
    % Right: Slide numbers
    \begin{beamercolorbox}[wd=.33\paperwidth,ht=2.25ex,dp=1ex,right]{date in head/foot}%
      \insertframenumber{} / \inserttotalframenumber\hspace*{2ex}
    \end{beamercolorbox}%
  }%
  \vskip0pt%
}

\usecolortheme{default}
\usefonttheme{serif}
\graphicspath{{../}{./}}
\setbeamertemplate{caption}[numbered]

% Title info
\title[Módulo 5: Tu Asistente Persona]{Módulo 5: Tu Asistente Personal 24/7}
\subtitle{Automatización y Finanzas Personales}
\author{Ricardo Reyes-Jimenez}
\institute{IA para Todos}
\date{\today}



\begin{document}

% Title slide
\begin{frame}
  \titlepage
\end{frame}

% Reset background for normal slides
\usebackgroundtemplate{}



\begin{frame}[fragile]{Objetivo de la Sesión}


\textbf{Al finalizar este módulo, los participantes serán capaces de:}

1. Comprender cómo delegar tareas repetitivas de la vida personal.
2. Usar la IA como consultor para organizar sus finanzas u objetivos personales.
3. Descubrir que el límite de la IA es la propia imaginación y necesidad de cada uno.

\end{frame}



\begin{frame}[fragile]{1. Ordenando las Finanzas Personales}


\textbf{Una hoja de cálculo puede ser aburrida; un asesor de finanzas con IA no.}

\begin{itemize}
  \item \textbf{Caso Práctico:} Quieres ahorrar para comprar un coche en dos años, pero no sabes cómo organizar tus gastos mensuales.
  \item \textbf{Uso de IA:} ``Actúa como un asesor financiero personal. Gano X al mes y mis gastos fijos son Y. Quiero ahorrar para un enganche de 5,000 dólares en 24 meses. Analiza mi capacidad de ahorro y dame un plan mensual realista, incluyendo categorías de dónde podría recortar gastos superfluos.''
\end{itemize}
\textbf{Resultado:} Un plan estructurado y personalizado sin necesidad de ser un experto contable.

\end{frame}



\begin{frame}[fragile]{2. El Nutriólogo y Entrenador en el Bolsillo}


\textbf{La IA te ayuda a mantener el rumbo hacia tus objetivos de salud.}

\begin{itemize}
  \item \textbf{Caso Práctico:} Tienes ingredientes específicos en la nevera y necesitas recetas saludables para la semana.
  \item \textbf{Uso de IA:} ``En mi refrigerador tengo pollo, espinacas, huevos, arroz y tomates. Quiero un menú de almuerzos para los próximos 3 días que sea bajo en carbohidratos. Dame la lista de recetas, las calorías aproximadas y los pasos de preparación.''
\end{itemize}
\textbf{Resultado:} Menos desperdicio de comida y una alimentación más alineada a tus metas.

\end{frame}



\begin{frame}[fragile]{3. Automatizando lo Repetitivo}


\textbf{Si haces una tarea digital más de tres veces a la semana, probablemente se pueda automatizar.}

\begin{itemize}
  \item \textbf{Caso Práctico:} Gestionar correos, recordatorios y citas puede consumir mucho tiempo de un freelancer o comerciante.
  \item \textbf{Soluciones Avanzadas (La punta del iceberg):} Herramientas como Zapier integradas con IA pueden leer correos de clientes, agendar citas automáticamente en tu calendario de Google y enviar confirmaciones, todo mientras tú te enfocas en el trabajo real.
\end{itemize}
\textbf{Reflexión:} La IA no solo crea textos, también actúa como el ``pegamento'' entre tus aplicaciones diarias.

\end{frame}



\begin{frame}[fragile]{Recapitulando el Módulo}


\textbf{Hoy hemos aprendido:}
\begin{itemize}
  \item Que la IA puede organizar datos cotidianos como presupuestos y menús semanales.
  \item A usar la IA como un asesor personal 24/7.
  \item Que el límite de uso lo define tu propia curiosidad y necesidad de ahorrar tiempo.
\end{itemize}
\textbf{¡Sigan aprendiendo, sigan trabajando y sigan disfrutando de la Inteligencia Artificial!}
\end{frame}




\begin{frame}[allowframebreaks]{References}
  \bibliographystyle{apalike}
  \bibliography{references}
\end{frame}


\end{document}